\section{The conformal factor}\label{sec:elliptic}

The sources begin at radius $L$.  This separation makes their effect near
$S$ of order $L^{-2}$, while their total integral is of order $L^{-1}$.
We prove both estimates by constructing the solution directly.

\begin{proposition}[Uniform conformal factor]\label{prop:conformal-factor}
There is $L_0>4$, depending only on the fixed cutoff, such that for every
$L\ge L_0$ and $s\in[0,1]$, Equation~\eqref{eq:scalar} has a smooth
axisymmetric solution $u=u_0+v$ on $\Omega$, where $u_0=1+r^{-1}$,
satisfying
\begin{equation}\label{eq:robin}
 u_r+\tfrac12u=0\quad\text{on }S,
 \qquad u\longrightarrow1\quad\text{as }r\longrightarrow\infty.
\end{equation}
It is unique among smooth solutions with $v\ge0$ and $v=O(r^{-1})$.
One may take a constant $M=3$, independent of $s,L$, such that
$1\le u_0\le u\le M$.  Uniformly in these parameters,
\begin{equation}\label{eq:radial-bounds}
 \begin{aligned}
  |v|+r|v_r|&\le Cs^2L^{-2} &&(r\ge1),\\
  |v_{rr}|&\le Cs^2L^{-2} &&(1\le r\le2).
 \end{aligned}
\end{equation}
For every integer $k\ge0$, all Cartesian derivatives satisfy
\begin{equation}\label{eq:all-derivative-bounds}
 |D^kv|\le C_k\frac{s^2}{L}r^{-1-k},
 \qquad |D^k(u-1)|\le C_k r^{-1-k}.
\end{equation}
For fixed $L$, $s\mapsto v_s$ is continuous in the sup norm, as are
\[
 s\longmapsto\int_{\R^3}H_s\,\dd x,
 \qquad
 s\longmapsto N_s(0)=\frac1{4\pi}\int_{\R^3}\frac{H_s(x)}{|x|}\,\dd x,
\]
where $H_s=pu_s^{-7}+\ell u_s^{-3}$ is extended by zero inside the unit
ball and $N_s$ is its Newton potential.  Moreover,
\begin{equation}\label{eq:source-integrals}
 0\le\int_{\R^3}H_s\,\dd x\le C\frac{s^2}{L},
 \qquad 0\le N_s(0)\le C\frac{s^2}{L^2}.
\end{equation}
At $s=0$ the solution is exactly $u_0$.
\end{proposition}

Here and below $D^k$ denotes any Cartesian derivative of order $k$;
constants may depend on this order but are independent of
$s\in[0,1]$ and $L\ge L_0$.
We first record the elementary potential estimates used in the proof.

\begin{lemma}[Newton potential and Robin reflection]\label{lem:newton}
Suppose $L>4$ and $H\in C(\R^3)$ obeys
\[
 |H(y)|\le a|y|^{-4}\mathbf1_{\{|y|\ge L\}}.
\]
Then
\[
 N_H(x)=\frac1{4\pi}\int_{\R^3}\frac{H(y)}{|x-y|}\,\dd y
\]
is well defined and $C^1$, satisfies $-\Delta N_H=H$ distributionally,
and is smooth and harmonic on $\{|x|<L\}$.  For $j=0,1$,
\begin{equation}\label{eq:newton-bounds}
 |D^jN_H(x)|\le
 \begin{cases}
  C_j aL^{-2-j},& |x|\le L/2,\\
  C_j a\bigl(r^{-2-j}+L^{-1}r^{-1-j}\bigr),&r=|x|\ge L/2.
 \end{cases}
\end{equation}
The first estimate holds for every $j\ge0$.  The operator
\[
 (\mathcal RH)(x)=N_H(x)+r^{-1}N_H(x/r^2),\qquad x\in\Omega,
\]
preserves nonnegativity, satisfies
$-\Delta\mathcal RH=H$ distributionally and
$(\partial_r+\tfrac12)\mathcal RH=0$ on $S$, and obeys
\begin{equation}\label{eq:robin-operator-bounds}
 \begin{gathered}
 \|\mathcal RH\|_\infty\le CaL^{-2},\qquad
 |\mathcal RH|+r|\partial_r\mathcal RH|\le CaL^{-2},\\
 |\partial_r^2\mathcal RH|\le CaL^{-2}\quad(1\le r\le2).
 \end{gathered}
\end{equation}
The homogeneous harmonic Robin problem on $\Omega$ has no nonzero
solution of order $O(r^{-1})$.
\end{lemma}

\begin{proof}
The bound on $H$ gives $\|H\|_{L^1}\le4\pi a/L$.  The Newton kernel
and its first derivatives are locally integrable.  Localizing around a
point and separating a small ball about its kernel singularity proves
that the potential and its first derivatives are continuous; the
singular part of the gradient integral is bounded by a constant times
the ball's radius.  The fundamental-solution identity
$-\Delta(4\pi|x|)^{-1}=\delta_0$, obtained by integration by parts around
the pole, therefore gives the distributional equation by Fubini's
theorem.

For $|x|\le L/2$, source separation allows differentiation under the
integral to every order and gives
\[
 |D^jN_H(x)|\le C_j a\int_L^\infty
          \rho^{-4}\rho^{-1-j}\rho^2\,\dd\rho
       \le C_j aL^{-2-j}.
\]
For $r=|x|\ge L/2$, split the integral at $|x-y|=r/2$.
In the inner part $|y|\ge r/2$, and for $j=0,1$ its absolute value is at
most
\[
 Ca r^{-4}\int_{|z|\le r/2}|z|^{-1-j}\,\dd z
 \le Ca r^{-2-j}.
\]
The outer part is bounded by
$Cr^{-1-j}\|H\|_{L^1}\le CaL^{-1}r^{-1-j}$.
This proves \eqref{eq:newton-bounds}.

Write $I(x)=x/r^2$.  The Kelvin identity in dimension three is
\[
 \Delta\bigl(r^{-1}N_H(I(x))\bigr)
       =r^{-5}(\Delta N_H)(I(x)).
\]
Its right side vanishes on $\Omega$ because $|I(x)|\le1<L$.
Along each ray $x=r\omega$, direct differentiation gives
\[
 (\mathcal RH)(\omega)=2N_H(\omega),\qquad
 \partial_r(\mathcal RH)(\omega)=-N_H(\omega),
\]
which proves the Robin condition including its angular dependence.
Positivity follows from the two positive kernels.  Combining
\eqref{eq:newton-bounds} in the two radial regions gives the first two
bounds in \eqref{eq:robin-operator-bounds}.  For the last bound both
potential arguments are separated from the source.  More explicitly,
put $n(t)=N_H(t\omega)$; then
\[
 \frac{\dd^2}{\dd r^2}\bigl(r^{-1}n(r^{-1})\bigr)
 =2r^{-3}n(r^{-1})+4r^{-4}n'(r^{-1})+r^{-5}n''(r^{-1}).
\]
The separated estimates for $j=0,1,2$, together with the corresponding
bound for $\partial_r^2N_H(r\omega)$, give the claim.

For uniqueness of the homogeneous problem, let $w$ be harmonic on
$\Omega$, $C^1$ up to $S$, with $w=O(r^{-1})$ and
$w_r+\tfrac12w=0$ on $S$.  Extend it to $0<r<1$ by
$r^{-1}w(I(x))$.  Values and tangential derivatives match at $S$, and
the interior radial derivative there is $-w-w_r=w_r$.
Thus the extension $\widehat w$ is distributionally harmonic across $S$
and hence harmonic on $\R^3\setminus\{0\}$. It is bounded near zero
and is $O(r^{-1})$ at infinity. For sufficiently small $\varepsilon$ and
large $R$, the maximum principle applied to $\pm\widehat w$ on
$\varepsilon<|x|<R$ gives
$|\widehat w(x)|\le C\varepsilon/|x|+C/R$, with one constant $C$:
the harmonic right side bounds the absolute value on both boundary
spheres. Letting $\varepsilon\downarrow0$ and $R\to\infty$ proves
uniqueness.
\end{proof}

\begin{lemma}[Localized derivative bootstrap]\label{lem:regularity}
Let $H$ satisfy the hypotheses of \Cref{lem:newton}.  If $H\in C^j(\R^3)$,
then $N_H\in C^{j+1}(\R^3)$.  If, in addition,
\[
 |D^iH(y)|\le a_j|y|^{-4-i}\mathbf1_{\{|y|\ge L\}}
 \quad(0\le i\le j),
\]
then \eqref{eq:newton-bounds} holds through order $j+1$, with a constant
$C_j a_j$ in place of $C_j a$.  Consequently the local gain in regularity
and the derivative decay use no first moment of $H$.
\end{lemma}

\begin{proof}
Fix a point $x_0$.  Choose a smooth compactly supported cutoff $\zeta$
equal to one near $x_0$, and keep this cutoff fixed when differentiating
with respect to $x$.  Writing $G(x)=(4\pi|x|)^{-1}$, for
$|\alpha|=j+1$ choose $\alpha=\beta+e_i$, $|\beta|=j$.  Distributional
integration by parts gives, near $x_0$,
\begin{equation}\label{eq:localized-kernel}
 D^\alpha N_H(x)
 =\int_{\R^3}\partial_iG(x-y)D_y^\beta(\zeta H)(y)\,\dd y
  +\int_{\R^3}D^\alpha G(x-y)(1-\zeta(y))H(y)\,\dd y.
\end{equation}
There are no boundary terms since $\zeta H$ has compact support.
The first integral is continuous by local integrability of
$\partial_iG$ and the same small-ball argument as above; the second is
smooth near $x_0$ by source separation and integrability at infinity.
The formula, also applied at lower orders, proves $C^{j+1}$ regularity.

For the estimate at $|x_0|=r\ge L/2$, take $\zeta$ supported in
$B(x_0,r/2)$ and equal to one on $B(x_0,r/4)$, with
$|D^i\zeta|\le C_i r^{-i}$.  On its support every Leibniz term in
$D^\beta(\zeta H)$ is bounded by $C_j a_jr^{-4-j}$.
The first integral in \eqref{eq:localized-kernel} at $x_0$ is therefore
bounded by
\[
 C_j a_jr^{-4-j}\int_{|z|\le r/2}|z|^{-2}\,\dd z
 \le C_j a_jr^{-3-j}.
\]
The second is bounded by
\[
 C_jr^{-2-j}\|H\|_{L^1}
 \le C_j a_jL^{-1}r^{-2-j}.
\]
These are precisely the claimed bounds at order $j+1$; the same
argument gives the lower orders.  On $|x|\le L/2$, differentiation of
the separated integral already gives every derivative estimate.
\end{proof}

\begin{proof}[Proof of \Cref{prop:conformal-factor}]
The coefficient estimates from \Cref{lem:seeds,lem:conformal} imply,
for every $j\ge0$,
\begin{equation}\label{eq:elliptic-coefficient-bounds}
 |D^jp|\le C_js^4r^{-6-j}\mathbf1_{\{r\ge L\}},\qquad
 |D^j\ell|\le C_js^2r^{-4-j}\mathbf1_{\{r\ge L\}}.
\end{equation}
Let $\mathcal C=\{v\in C_b(\Omega):v\ge0\}$, a complete metric space
in the sup norm.  For $v\in\mathcal C$, put
\begin{equation}\label{eq:fixed-point}
 \begin{split}
 H_v&=p(u_0+v)^{-7}+\ell(u_0+v)^{-3},\\
 N_v(x)&=\frac1{4\pi}\int_{\R^3}\frac{H_v(y)}{|x-y|}\,\dd y,\\
 T_s(v)(x)&=N_v(x)+r^{-1}N_v(x/r^2).
 \end{split}
\end{equation}
The source is extended by zero to $\R^3$; this extension is continuous
because it vanishes on $r<L$.  Since $u_0+v\ge1$, the mean-value theorem
for the negative powers gives
\begin{equation}\label{eq:source-contraction}
 \begin{split}
 0\le H_v&\le Cs^2r^{-4}\mathbf1_{\{r\ge L\}},\\
 |H_v-H_w|&\le Cs^2r^{-4}\mathbf1_{\{r\ge L\}}
                  \|v-w\|_\infty.
 \end{split}
\end{equation}
By \Cref{lem:newton}, $T_s$ maps the whole cone $\mathcal C$ into itself,
with
\[
 \|T_s(v)\|_\infty\le Cs^2L^{-2},\qquad
 \|T_s(v)-T_s(w)\|_\infty\le CL^{-2}\|v-w\|_\infty.
\]
Choose $L_0>4$ so large that the second constant is at most $1/2$ and
the first, for $s=1$, is at most one.  The contraction theorem now gives
an actual fixed point $v\in\mathcal C$, unique in that cone.
The operator commutes with rotation about the prescribed axis, so its
unique fixed point is axisymmetric.  It also gives $v=0$ when $s=0$.
\Cref{lem:newton} yields \eqref{eq:radial-bounds}, $1\le u_0\le u\le3$,
the Robin condition, and the scalar equation distributionally.
The potential bounds give $v=O(r^{-1})$, hence $u\to1$.

For completeness we bootstrap both smoothness and the full decay.
The fixed point is $C^1$ by \Cref{lem:newton}; near $S$ it is already
smooth by separation of both potential arguments from the source.
Suppose derivatives of $u-1$ through order $j$ have the uniform bounds
$C_i r^{-1-i}$.  The case $j=0$ follows from
\eqref{eq:newton-bounds} and $u_0=1+r^{-1}$.  Because $u\ge1$, the
chain and product rules with \eqref{eq:elliptic-coefficient-bounds}
give
\[
 |D^iH_v|\le C_js^2r^{-4-i}\mathbf1_{\{r\ge L\}}
 \quad(0\le i\le j).
\]
The zero extension is $C^j$ since the source vanishes near $S$.
\Cref{lem:regularity} gains one derivative for $N_v$ with the stated
bounds.  For the reflected term, the separated estimates and the
homogeneity of $I(x)$ give at every order $k$
\[
 \bigl|D^k\bigl(r^{-1}N_v(I(x))\bigr)\bigr|
       \le C_ks^2L^{-2}r^{-1-k}.
\]
In the region $1\le r\le L/2$, the ordinary term is bounded by
$C_ks^2L^{-2-k}$; in the other region it is bounded by
$C_ks^2(r^{-2-k}+L^{-1}r^{-1-k})$.
Each bound is at most $C_ks^2L^{-1}r^{-1-k}$ in its region.
This proves the next induction step and \eqref{eq:all-derivative-bounds}
at every order, with constants independent of $L,s$.
In particular $u$ is smooth and solves \eqref{eq:scalar} pointwise.
If another solution has $v\ge0$ and $v=O(r^{-1})$, subtracting
$\mathcal RH_v$ leaves a decaying homogeneous harmonic Robin solution.
\Cref{lem:newton} makes that difference zero; cone uniqueness then
identifies the solutions.

Finally, fix $L\ge L_0$.  The coefficients have the form
$p_s=s^4p_1$, $\ell_s=s^2\ell_1$.  Thus uniformly for $v\in\mathcal C$,
\[
 \|T_s(v)-T_t(v)\|_\infty\le CL^{-2}|s-t|.
\]
Using the contraction constant $1/2$ in the fixed point equations gives
\[
 \|v_s-v_t\|_\infty\le 2CL^{-2}|s-t|.
\]
Combining this with the pointwise source Lipschitz bound yields
\[
 |H_s(y)-H_t(y)|\le C|s-t|\,|y|^{-4}
                    \mathbf1_{\{|y|\ge L\}}.
\]
The right side is integrable with both weights $1$ and $|y|^{-1}$.
It proves continuity of the two source integrals, and
\eqref{eq:source-contraction} gives their quantitative bounds
\eqref{eq:source-integrals}.  Uniform convergence also gives continuity
of the restriction $u_s|_S$, which will imply continuity of the horizon
area.  No finite first moment or multipole expansion has entered the
argument.
\end{proof}
